% 06_continuous_population.tex

\section{Continuous-Time Population Dynamics}
\label{sec:continuous-population}

Section~\ref{sec:belief-loop} gives the strategy-level statement: for \(q>q^*\), no D1-consistent pure assessment has stable top-end occupancy of the ambitious signal. After Section~\ref{sec:polarization} records the cross-sectional shadow, this section gives the population/time-level counterpart. It treats ambitious participation as a continuous population stock and asks what the cross section looks like once the same D1-closed payoff map is imposed.

The correspondence is:
\[
\text{no occupied ambitious fixed point}
\quad\Longleftrightarrow\quad
\text{demographic exit from the top-end ambitious pool}.
\]
The continuous model is not a chronological avalanche story. It is an economic packaging of the same closed-loop state claim: D1 closure turns top-end ambitious participation into an exit process, and the stationary boundary profile records where the remaining ambitious density can be sustained.

\subsection{Reaction-Diffusion Occupancy}

Let \(x\in[0,\infty)\) denote a continuous economic state. Higher \(x\) means higher visible status, stronger reputation-at-risk, or more self-accurate exposure to future adverse repricing. Let \(n(x,t)\ge0\) be the mass density of agents currently occupying the ambitious signal \(s_A\) at state \(x\). The relevant object is the D1-closed aggregate ambitious pool, not separate type-by-type participation.

Local mobility in status space is summarized by a diffusion term. Local entry, retention, and exit from the ambitious pool are summarized by the D1-closed growth rate
\[
V_q(x)=B^{D1}(x)-\Gamma(q,\eta)\Phi(x).
\]
Here \(B^{D1}(x)\) is the current net return to ambitious occupancy after the off-path D1 map has selected the low-attribution assessment, while \(\Gamma(q,\eta)\Phi(x)\) is the continuation exposure price inherited from Section~\ref{sec:belief-loop}. The function \(\Phi(x)\) is increasing enough near the top that
\[
V_q(x)<0
\]
on the exposed upper region. In that region the demographic exit rate from the ambitious pool is
\[
\kappa_q(x):=-V_q(x)>0.
\]

The resulting continuous-time occupancy equation is
\begin{equation}
\frac{\partial n}{\partial t}
=
\frac{\sigma^2}{2}\frac{\partial^2 n}{\partial x^2}
+V_q(x)n.
\label{eq:occupancy-pde}
\end{equation}
The first term lets ambitious density diffuse across nearby economic states. The second term is birth-death accounting inside the ambitious pool: \(V_q(x)>0\) permits local replenishment, while \(V_q(x)<0\) kills ambitious occupancy at rate \(\kappa_q(x)\). The D1 closure matters because it changes the top-end payoff from a reputational rent into a negative participation margin.

\subsection{Long-Run Cross-Sectional Density}

A stable long-run cross-sectional density has the form
\[
n(x,t)=e^{\lambda t}\psi(x),
\qquad
\psi(x)\ge0,
\]
with normalization supplied by the protected lower region. Substituting this form into equation~\eqref{eq:occupancy-pde} gives
\begin{equation}
\frac{\sigma^2}{2}\psi''(x)+V_q(x)\psi(x)=\lambda\psi(x).
\label{eq:stationary-density-eigen}
\end{equation}
The scalar \(\lambda\) is the net long-run growth rate of the ambitious stock under the D1-closed occupancy process. When the remaining density is normalized to unit mass, \(\psi\) is the stable long-run cross-sectional shape.

The boundary conditions have direct economic meanings. At the protected lower limit, the model either anchors base density or fixes a normalization:
\[
\psi(0)=\psi_0>0
\qquad\text{or}\qquad
\int_0^\infty \psi(x)\,dx=1.
\]
At the far top end, the non-explosive transversality condition rules out profiles that create infinite ambitious mass:
\[
\limsup_{x\to\infty}\psi(x)<\infty,
\qquad
\psi\in L^1_+(0,\infty).
\]
If the upper tail has \(V_q(x)\le -\underline\kappa<0\), this condition selects a decaying boundary profile. Population density cannot be negative, and an admissible long-run cross section cannot alternate signs. Oscillatory continuations are therefore not economic population densities.

\subsection{Critical Boundary Profile}

Let \(\bar x\) be the critical top boundary where the D1-closed growth rate changes sign relative to the long-run growth rate. Near that boundary, use the turning-point approximation
\[
V_q(x)\approx a(q)-b(\bar x-x),
\qquad b>0.
\]
With the inward coordinate
\[
y=\bar x-x,
\]
equation~\eqref{eq:stationary-density-eigen} becomes
\begin{equation}
\frac{\sigma^2}{2}\psi''(y)+\bigl(a(q)-by\bigr)\psi(y)=\lambda\psi(y).
\label{eq:turning-point-localized}
\end{equation}
After the rescaling
\[
z=
\left(\frac{2b}{\sigma^2}\right)^{1/3}
\left(
y-\frac{a(q)-\lambda}{b}
\right),
\]
the local equation takes the standard turning-point form
\[
\frac{d^2\psi}{dz^2}-z\psi=0.
\]
Mathematically, this is the canonical Airy turning-point solution. Economically, only the nonnegative bounded branch is admissible: the growing branch violates non-explosive transversality, and the oscillatory continuation would imply negative population density. The relevant object is the nonnegative lobe on the protected side together with the decaying boundary-layer tail on the top side.

\begin{figure}[ht]
    \centering
    \includegraphics[width=0.85\textwidth]{airy_density_corrected.pdf}
    \caption{Critical boundary profile. Near the cutoff where the D1-closed growth rate changes sign, the stationary occupancy equation has the standard turning-point form. The decaying boundary-layer profile represents depletion of ambitious occupancy at the top end.}
    \label{fig:critical-boundary-profile}
\end{figure}

Figure~\ref{fig:critical-boundary-profile} is a local critical snapshot, not a new global equilibrium theorem. The protected side represents states where ambitious occupancy can be sustained or replenished. The exposed top side represents states where \(V_q(x)<0\), so demographic exit dominates local diffusion and the density thins into a depletion tail.

\subsection{Economic Implication}

The continuous population formulation adds one distributional implication to the fixed-point theorem. D1 closure removes the high-type rent from ambitious occupancy; at sufficiently exposed top states, the remaining D1-closed payoff is negative. The reaction-diffusion process then has a top-end depletion layer rather than a refilled ambitious upper tail.

Thus the continuous model is the macroscopic shadow of the belief-loop fixed point. Strategy level: no D1-consistent fixed point has stable top-end occupancy of \(s_A\). Population level: once \(V_q(x)<0\) near the critical top boundary, ambitious density is killed in that region and the admissible long-run profile decays. Continuous population diffusion does not refill top-end ambitious occupancy once the D1-closed exit rate dominates.
